% Kahlus one-pager for Amrith Lotlikar (Stanford EE).
% Source of truth: Obsidian "Amrith one page.md" (corrected 2026-08-19)
%   + artifacts/isolated_residual_evidence_20260815/VERDICT.md (added 2026-08-20)
% Build: pdflatex main.tex   (no bibliography, no shell-escape, Overleaf-safe)
\documentclass[11pt]{article}

\usepackage[letterpaper,margin=1in]{geometry}
\usepackage[T1]{fontenc}
\usepackage[utf8]{inputenc}
\usepackage{charter}
\usepackage[expansion=false]{microtype}
\usepackage{booktabs}
\usepackage{xcolor}
\usepackage{titlesec}
\usepackage[hidelinks]{hyperref}
\usepackage{url}
\urlstyle{tt}

\definecolor{navy}{HTML}{173B63}
\definecolor{graytext}{HTML}{555555}

\titleformat{\section}{\normalfont\bfseries\color{navy}}{}{0pt}{}
\titlespacing*{\section}{0pt}{10pt}{4pt}

\setlength{\parindent}{0pt}
\setlength{\parskip}{7pt}
\pagestyle{empty}

\newcommand{\lead}[1]{\textbf{\color{navy}#1}}

\begin{document}

{\LARGE\bfseries Kahlus: what the retraction bought}\\[2pt]
{\color{graytext}\small Aayushya Patel \quad|\quad for Amrith Lotlikar \quad|\quad 21 August 2026}

\vspace{4pt}
\hrule
\vspace{8pt}

\begin{center}
\fcolorbox{navy}{yellow!15}{%
\parbox{0.96\linewidth}{\footnotesize
\textbf{For Amrith.} Tests start 22 August 2026. Everything below this box is
the current evidence only. None of it is a capability claim for a model we
have not yet run.

\smallskip
\textbf{Build plan} (the levers that make the model better, not the story),
ranked by expected gain:

1. Stop fighting where the war is over. At 10--160\,ms the linear model owns
almost all extractable signal. Headroom is at $320$\,ms and beyond: the stack
sits at MSE $\approx 0.80$ against a mean-reversion wall near $1.0$, and nobody
knows the true floor out there. Every lever below attacks that region.

2. Give the model eyes. We feed it 1 second, one channel. Sleep EEG runs on
quasi-periodic machinery (spindles, K-complexes, slow waves) at 10--30\,s.
Sleep-EDF also has a second EEG channel plus EOG and EMG, unused. Longer
context and multi-channel is a harness extension, not a research bet. The
sweep script is already in the repo. About $70\%$ this beats today's stack
visibly at long horizons.

3. Build the architecture the physics asks for. The oscillator rung is dumb
fixed FFT bins. Right inductive bias for a quasi-periodic signal: a learned
resonant state-space (S4-family, poles near the unit circle). About
$30$--$40\%$ it beats the GRU stack. If it does, that is a small named model
class.

4. Scale data $13\times$ where scaling might actually pay. Challenge-2018
($\sim$1{,}000 subjects) on the cluster. If the nonlinear residual grows with
data, we ride it. If it flattens, we stop wasting GPUs where physics says no.

5. Change the target where microvolts run out. Sleep-onset in the next $N$
seconds is Track 3 of the EEG/EMG Foundation Challenge 2026 attached to
BrainBodyFM (1{,}197 teams last year, winner gets a spotlight). That is the
public scoreboard.

\smallskip
\textbf{Sequence.} Paper ships 5 Sep 2026 on \emph{current} results, unchanged.
In parallel: \#2 now (capability curve by the weekend), then \#3, then \#4
after Krish's current job, then \#5 after 5 Sep.

\textbf{One thing only you can do:} register us for the EEG/EMG Foundation
Challenge 2026.

\textbf{Tonight.} Both. Build now. Launch the moment the GPU frees.
}}
\end{center}

\vspace{6pt}

\lead{The problem.}
EEG forecasting papers report scores that look close to solved. Ours did too:
MSE 3.116 and $r = 0.972$ from a 100{,}000-step recurrent model on
BNCI2014\_001, against ridge at 7.745 (the forecast sidecar of a multitask
model whose aggregate score was MSE 28.55, $r = 0.48$). That number is not
forecasting. The recovered target shared 126 of its 127 samples with the
input window, so a model
that copies its own input scores well. For squared error the inflation has a
closed form, $W/H$, where $W$ is the scored window and $H$ the horizon. We
checked 286 configurations and found no violations, and the bound predicts the
$126\times$ gap we actually measured. % NOTE: correlation-law numbers (223 cells, s_pred~0.70, r=0.984 vs 0.014)
% are the F10 artifact on PR #77; not yet merged to main.
Correlation is worse and does not obey
that shortcut: the exact factor is
$1 + \frac{W-H}{H}\left(s_{\mathrm{pred}} s_{\mathrm{targ}}\right)^{-1/2}$,
and because models shrink toward the mean where signal is absent
($s_{\mathrm{pred}} \approx 0.70$), $W/H$ \emph{understates} the true inflation
on 223 of the same 286 cells, always in the flattering direction. The sharpest
case is white noise at $W{=}32$, $H{=}1$ reporting $r = 0.984$ when the true
forecast correlation is $0.014$. So $0.972$ is retracted for a stronger reason
than $3.116$ is.

\lead{What the field does, and why it does not settle this.}
Deep EEG models are usually compared against weak baselines, or trained and
tested on windows from the same person. The 2026 foundation-model benchmarks
make the same point for classification: tuned specialists stay competitive,
linear probes are weak, and scale is not a free win. Forecasting shows up only
as a side task. FoME forecasts patches on a 6:2:2 split and reports no
persistence or ridge baseline. So there is no standard test where a model must
beat copy-the-last-sample and a linear map on a strictly future sample from a
subject it has never seen.

\lead{What we did, and where it stands.}
We built that test and ran our own model through it. Splits are frozen before
windowing, subjects are held out, and the scored sample is never in the input.
Baselines are persistence, ridge, and mean. On Sleep-EDF the GRU beats both at
every horizon we ran except 64 samples ahead: 0.211 against ridge 0.227 and
persistence 0.281 at one sample, 0.690 against 0.822 and 1.334 at sixteen. That
margin is not established as significant, and on 9-subject BNCI LOSO the same
protocol cannot separate the GRU from tuned ridge at any horizon --- 0 of 9
horizons reach $p < 0.05$. The contribution is the protocol and the retraction,
not a new model.

\section{Since we met: the residual run came back positive, and I do not trust it yet}

At our meeting this was the next experiment. It ran on 15 August, on Sleep-EDF
78-subject LOSO: strip what ridge can predict, strip narrowband oscillations,
ask whether a GRU finds anything in the leftover.

It won at every horizon. Because an undertuned linear baseline has already
manufactured one win in this project, the first thing we did was attack it:
re-score against a 127-lag ridge with $\alpha$ tuned per fold and per horizon
on held-out training subjects, identical windows, identical train-fold-only
normalisation. Tuning the baseline closed 5--49\% of the raw gap.

\begin{center}
\begin{tabular}{lccccc}
\toprule
horizon & weak ridge & tuned ridge & GRU stack & gap [95\% CI] & folds won \\
\midrule
1 (10 ms)   & 0.2951 & 0.2932 & 0.2600 & $+0.0332$ $[+0.0280, +0.0392]$ & 77/78 \\
8 (80 ms)   & 0.6603 & 0.6370 & 0.5418 & $+0.0952$ $[+0.0799, +0.1119]$ & 78/78 \\
16 (160 ms) & 0.9004 & 0.8410 & 0.6929 & $+0.1481$ $[+0.1219, +0.1774]$ & 75/78 \\
32 (320 ms) & 0.9872 & 0.8970 & 0.8025 & $+0.0945$ $[+0.0720, +0.1239]$ & 72/78 \\
\bottomrule
\end{tabular}
\end{center}

A 10.5--17.6\% MSE advantage survives a tuned baseline at every horizon,
per-horizon CIs excluding zero, better on 72--78 of 78 held-out subjects. I am calling it a
\emph{powered preliminary positive}, not a result, because five things are open:
one seed where our own standard is three; an undisclosed normalisation change in
the run script that I have not reproduced from repo code; one channel of one
corpus; BNCI still negative under a weaker fixed-split protocol; and the
baseline ladder stops at ridge, where the gate demands the best available
baseline, so gradient boosting is next.

There is also a direct tension I would rather show you than smooth over. F8 put
the plain-GRU advantage at zero by about 320 ms. This stack still shows
$+0.0945$ at 320 ms on 72 of 78 subjects. Different architecture, different
split, different normalisation --- so it is not a contradiction yet, but it is
the kind of gap that resolves into either a real finding or a bug, and I would
like your read on which controlled comparison settles it fastest.

\section{The two questions I actually want answered}

\textbf{1. Is the residual stack worth a paper on its own, or is it Result 3 of
the protocol paper?} My instinct is the latter: the protocol plus the
retraction is the contribution, and a nonlinear residual that survives a tuned
baseline is the evidence that the protocol is worth adopting. But that may be
undervaluing the only positive result in the project.

\textbf{2. What kills this fastest?} The seed sweep and the GBM baseline are
both cheap. If you think either is likely to erase an 11--18\% margin, I would
rather run it before writing anything than after.

\vspace{6pt}
\hrule
\vspace{4pt}
{\color{graytext}\small
Every number above is read from a run artifact, not from memory.

\smallskip
Overlap work (public): \url{github.com/ptlnextdoor/eeg-forecastability-limits}

Residual pack: \path{artifacts/isolated_residual_evidence_20260815/}
--- per-fold scores, tuned-ridge audit, SHA-verified handoff.

\smallskip
Seven self-caught corrections are listed in the handbook's honest ledger,
including two I fixed in this page on 19 August.}

\end{document}
